\chapter{مبانی و تعریف مسئله}
\label{ch:preliminaries}

\section{دسته‌بندی و تابع هزینه}
فرض می‌کنیم $x\in[0,1]^d$ یک تصویر با $d$ مؤلفه و $y\in\{1,\ldots,C\}$ برچسب آن باشد. شبکه $f_\theta:\IR^d\to\IR^C$ با پارامتر $\theta$، بردار امتیازهای پیش از نرمال‌سازی را تولید می‌کند. دسته پیش‌بینی‌شده برابر است با
\begin{equation}
\widehat y_\theta(x)=\mathop{\arg\max}_{c\in\{1,\ldots,C\}} f_\theta(x)_c.
\end{equation}
برای تبدیل امتیازها به توزیع احتمال از بیشینه هموار استفاده می‌شود:
\begin{equation}
q_\theta(c\mid x)=\frac{\exp(f_\theta(x)_c)}{\sum_{j=1}^{C}\exp(f_\theta(x)_j)}.
\end{equation}
تابع هزینه آنتروپی متقاطع برای برچسب درست به شکل
\begin{equation}\label{eq:crossentropy}
\mathcal L(f_\theta(x),y)=-\log q_\theta(y\mid x)
\end{equation}
تعریف می‌شود. در ادامه، هنگامی که $\theta$ و $y$ ثابت‌اند، از نماد کوتاه $\ell(x)$ برای همین هزینه استفاده می‌کنیم. تمایز میان هزینه و درستی ضروری است: هزینه به اطمینان پیش‌بینی نیز حساس است، اما درستی تنها درست یا نادرست‌بودن دسته نهایی را می‌شمارد.

در آموزش معمول، امید هزینه روی توزیع داده $\mathcal D$ کمینه می‌شود. در عمل، این امید با میانگین روی داده آموزشی و سپس روی دسته‌های کوچک تقریب زده می‌شود. اگر $B$ یک دسته باشد، هزینه دسته برابر است با
\begin{equation}
\mathcal L_B(\theta)=\frac{1}{|B|}\sum_{(x,y)\in B}\mathcal L(f_\theta(x),y).
\end{equation}
یک دوره آموزش، معمولاً یک پیمایش داده آموزشی است؛ یک به‌روزرسانی دسته تنها بخشی از آن است. ازاین‌رو، زمان هشدار برحسب دسته را نمی‌توان بدون دانستن تعداد دسته‌ها به دوره تبدیل کرد.

\section{مدل تهدید و مجموعه آشفتگی}
\begin{تعریف}[آشفتگی مجاز]
برای تصویر $x$ و بودجه $\epsilon>0$، مجموعه تهدید نرم بی‌نهایت را به صورت
\begin{equation}\label{eq:threatset}
\Delta_\infty(x,\epsilon)=\{\delta\in\IR^d:\|\delta\|_\infty\leq\epsilon,\quad x+\delta\in[0,1]^d\}
\end{equation}
تعریف می‌کنیم که در آن $\|\delta\|_\infty=\max_j|\delta_j|$ است.
\end{تعریف}
این قید بیشینه تغییر هر مؤلفه را محدود می‌کند. برای تصویر با مقیاس $[0,1]$، بودجه $8/255$ به معنای حداکثر تغییر هشت سطح از مقیاس هشت‌بیتی در هر مؤلفه است. این تعریف، یک قرارداد ریاضی برای ارزیابی است و به‌تنهایی تضمین نمی‌کند که هر آشفتگی مجاز از نظر ادراکی نامحسوس باشد. اهمیت یک بودجه ثابت می‌تواند برای دو دادگان با بافت و تفکیک‌پذیری متفاوت، یکسان نباشد.

حمله بدون هدف می‌کوشد مدل را از دسته درست خارج کند و معمولاً هزینه برچسب درست را افزایش می‌دهد. حمله هدفمند، خروجی را به دسته‌ای مشخص هدایت می‌کند. تمرکز این پژوهش بر حالت بدون هدف است. در ارزیابی جعبه سفید، حمله به پارامترها و شیب مدل دسترسی دارد؛ بنابراین، افت عملکرد را نمی‌توان صرفاً به محدودبودن اطلاعات مهاجم نسبت داد.

باید میان «اندازه گام» و «بودجه تهدید» تفاوت گذاشت. بودجه، مجموعه مجاز ارزیابی را تعریف می‌کند، اما اندازه گام، حرکت الگوریتم حمله را کنترل می‌کند. یک حمله تکراری می‌تواند گامی بزرگ‌تر از بودجه بردارد و سپس به مجموعه مجاز تصویر شود. در مقابل، اگر تصویرشدن به کره حذف شود، دامنه ورودی‌های آموزشی می‌تواند از کره ارزیابی بزرگ‌تر باشد. این تفاوت در تفسیر \sora{} نقش دارد و در فصل روش تصریح می‌شود.

\section{حمله علامت شیب سریع}
با ثابت‌گرفتن پارامترها، بسط مرتبه اول هزینه در $x$ می‌دهد
\begin{equation}\label{eq:firstorder}
\ell(x+\delta)\approx\ell(x)+g^\top\delta,\qquad g=\nabla_x\ell(x).
\end{equation}
بیشینه جمله خطی تحت قید $\|\delta\|_\infty\leq\epsilon$ با انتخاب مؤلفه‌ای علامت شیب به دست می‌آید:
\begin{equation}\label{eq:fgsm}
\delta_{\mathrm{FGSM}}=\epsilon\,\sign(g),\qquad
x_{\mathrm{adv}}=\Pi_{[0,1]^d}\bigl(x+\delta_{\mathrm{FGSM}}\bigr).
\end{equation}
این روش، حمله علامت شیب سریع یا \lr{FGSM} نام دارد \cite{goodfellow2015explaining}. دلیل انتخاب علامت آن است که هر مؤلفه مستقل از مؤلفه‌های دیگر، سهم خود را در ضرب داخلی بیشینه می‌کند. اگر $g_j>0$ باشد، انتخاب $\delta_j=\epsilon$ و اگر $g_j<0$ باشد، انتخاب $\delta_j=-\epsilon$ مناسب است. مؤلفه صفر شیب در تقریب خطی جهت ترجیحی ندارد.

این استدلال برای تابع خطی دقیق است، اما تابع هزینه یک شبکه عمیق در یک ناحیه بزرگ الزاماً خطی نیست. حمله تک‌گامی تنها شیب نقطه آغاز را می‌بیند. اگر جهت شیب یا خمیدگی در طول مسیر به‌شدت تغییر کند، نقطه انتهایی گام لزوماً بدترین نقطه مسیر نیست. بنابراین، شکست \lr{FGSM} در یافتن نمونه نامطلوب نمی‌تواند به‌تنهایی نشانه ستبری مدل باشد.

آغاز تصادفی یکی از راه‌های افزایش تنوع است. ابتدا $\eta$ از یک توزیع روی همسایگی ورودی نمونه‌برداری می‌شود و شیب در $x+\eta$ محاسبه می‌شود. بااین‌حال، انتخاب دامنه نویز، اندازه گام و وجود یا نبود تصویرشدن به کره، روش‌های متفاوتی می‌سازد. همه روش‌های دارای آغاز تصادفی یک حمله واحد نیستند.

\section{حمله چندگامی با تصویرشدن}
در حمله \lr{PGD}، تخمین آشفتگی به‌طور تکراری به‌روزرسانی می‌شود \cite{madry2018towards}:
\begin{align}
\delta^{(0)}&\sim\mathcal U([-\epsilon,\epsilon]^d),\\
\delta^{(t+1)}&=\Pi_{\Delta_\infty(x,\epsilon)}\left(
\delta^{(t)}+a\,\sign\!\left(\nabla_x\mathcal L(f_\theta(x+\delta^{(t)}),y)\right)\right).
\label{eq:pgd}
\end{align}
در این رابطه، $a$ اندازه گام و $t$ شماره گام است. نام تاریخی روش به نزول شیب اشاره دارد، ولی برای افزایش هزینه در مسئله درونی، جهت صعود به کار می‌رود. در قید نرم بی‌نهایت، تصویرشدن را می‌توان با برش مؤلفه‌ای پیاده‌سازی کرد:
\begin{equation}
\delta_j\leftarrow\min\{\epsilon,1-x_j,\max\{-\epsilon,-x_j,\delta_j\}\}.
\end{equation}
برش به دامنه پیکسل و برش به بودجه آشفتگی دو قید مستقل هستند.

افزایش تعداد گام‌ها و استفاده از چند آغاز تصادفی می‌تواند احتمال یافتن نقاط دشوارتر را افزایش دهد، اما موفقیت حمله به تنظیم درست اندازه گام و تابع هدف نیز وابسته است. بنابراین، نام \lr{PGD-10} به‌تنهایی تمام پروتکل ارزیابی را مشخص نمی‌کند. بودجه، اندازه گام، تعداد بازآغازها، تعداد نمونه آزمون و شیوه انتخاب نقطه کنترل نیز برای مقایسه لازم‌اند.

\section{آموزش هم‌آورد به‌عنوان بهینه‌سازی تودرتو}
صورت اصلی مسئله مورد مطالعه چنین است:
\begin{equation}\label{eq:minmax}
\min_\theta\;\mathbb E_{(x,y)\sim\mathcal D}
\left[\max_{\delta\in\Delta_\infty(x,\epsilon)}\mathcal L(f_\theta(x+\delta),y)\right].
\end{equation}
در هر به‌روزرسانی، ابتدا برای پارامترهای فعلی نمونه هم‌آورد ساخته می‌شود و سپس هزینه همان نمونه نسبت به پارامترها کمینه می‌شود. در پیاده‌سازی رایج، مسیر تولید آشفتگی از نمودار مشتق‌گیری به‌روزرسانی وزن‌ها جدا می‌شود؛ به‌عبارت دیگر، به‌روزرسانی پارامتر معمولاً مشتق کامل همه گام‌های مولد حمله را محاسبه نمی‌کند.

اگر مسئله درونی ضعیف حل شود، مسئله بیرونی روی هزینه‌ای کمتر از بیشینه واقعی آموزش می‌بیند. این خطا صرفاً یک تقریب عددی ثابت نیست؛ مدل ممکن است در طول آموزش شکل تابع هزینه را طوری تغییر دهد که حل‌کننده درونی بیشتر دچار خطا شود. بیش‌برازش فاجعه‌بار نمونه برجسته این برهم‌کنش میان مدل و مولد نمونه آموزشی است.

در آموزش سریع، هدف کاهش هزینه حل درونی است. اگر ساخت حمله $K$ بار مشتق نسبت به ورودی بخواهد، بخش حمله تقریباً با $K$ رشد می‌کند؛ بااین‌حال، زمان کل دقیقاً $K$ برابر نمی‌شود، زیرا بارگذاری داده، گذر پیشرو، به‌روزرسانی پارامتر و ویژگی‌های سخت‌افزار نیز مؤثرند. مقایسه هزینه باید با اندازه‌گیری روی شرایط یکسان یا با شمارش روشن گذرهای محاسباتی انجام شود.

\section{درستی پاک و درستی هم‌آورد}
برای مجموعه آزمون $\{(x_i,y_i)\}_{i=1}^n$، درستی پاک برابر است با
\begin{equation}
\widehat A_{\mathrm{clean}}=\frac1n\sum_{i=1}^n
\mathbf1\{\widehat y_\theta(x_i)=y_i\}.
\end{equation}
درستی ستبر واقعی در مدل تهدید مشخص، مستلزم درست‌بودن پیش‌بینی برای تمام آشفتگی‌های مجاز است:
\begin{equation}\label{eq:true-robust}
\widehat A_{\mathrm{rob}}=\frac1n\sum_{i=1}^n
\mathbf1\{\forall\delta\in\Delta_\infty(x_i,\epsilon):
\widehat y_\theta(x_i+\delta)=y_i\}.
\end{equation}
در عمل، حمله تنها برخی نقاط را می‌جوید. درستی تجربی در برابر یک حمله، معمولاً کران بالایی بر درستی ستبر واقعی است، زیرا شکست حمله در یافتن خطا به معنای نبود خطا نیست. چون صفر در مجموعه آشفتگی حضور دارد، درستی ستبر واقعی از درستی پاک بیشتر نمی‌شود؛ ولی درستی یک حمله خاص می‌تواند از درستی پاک بیشتر شود، زیرا حمله ممکن است برخی نمونه‌های ابتدا نادرست را به دسته درست جابه‌جا کند.

\lr{AutoAttack} مجموعه‌ای از حملات مکمل را برای ارزیابی پایاتر گرد می‌آورد \cite{croce2020reliable}. هدف آن کاهش اتکا به یک حمله یا یک انتخاب حساس ابرپارامتر است. بااین‌حال، این ارزیابی نیز تجربی است. در این پایان‌نامه، هرگاه «درستی هم‌آورد» با نام حمله گزارش می‌شود، منظور درستی تجربی تحت همان پروتکل است و نه گواهی جامع ستبری.

\section{شیب، هسین و هندسه محلی}
برای تابع هزینه هموار، شیب $g=\nabla_x\ell(x)$ جهت بیشترین افزایش مرتبه اول در هندسه اقلیدسی را نشان می‌دهد و هسین $H=\nabla_x^2\ell(x)$ تغییر شیب در همسایگی را توصیف می‌کند. بسط مرتبه دوم چنین است:
\begin{equation}\label{eq:quadratic}
\ell(x+\delta)=\ell(x)+g^\top\delta+\frac{1}{2}\delta^\top H\delta+R(\delta).
\end{equation}
اگر هسین در همسایگی لیپشیتس باشد، باقی‌مانده هزینه از مرتبه $\|\delta\|^3$ است. روی مسیر ثابت $\delta=\alpha v$، جمله درجه دوم تنها به کمیت اسکالر $v^\top Hv$ وابسته می‌شود. این کمیت، خمیدگی جهتی است و با کل هسین یکسان نیست.

محاسبه و نگهداری صریح یک ماتریس $d\times d$ برای تصاویر بزرگ پرهزینه است. در مقابل، دو شیب در نقاط نزدیک حاوی اطلاعاتی درباره $Hv$ هستند:
\begin{equation}\label{eq:gradient-difference}
\nabla_x\ell(x+\alpha v)-\nabla_x\ell(x)\approx\alpha Hv.
\end{equation}
این رابطه اساس استفاده غیرمستقیم از اطلاعات مرتبه دوم است. مزیت آن زمانی بیشتر می‌شود که هر دو شیب، به دلایل دیگری در آموزش محاسبه شده باشند.

شباهت کسینوسی دو بردار ناصفر $a$ و $b$ برابر است با
\begin{equation}
\cos(a,b)=\frac{a^\top b}{\|a\|_2\|b\|_2}.
\end{equation}
این کمیت جهت را مقایسه می‌کند و نسبت به ضرب هر بردار در یک اسکالر مثبت ناورداست. در نتیجه، تغییر بزرگی شیب بدون تغییر جهت ممکن است در آن دیده نشود. این ویژگی هم برای تفسیر \pertalign{} مفید است و هم حد اطلاعاتی آن را نشان می‌دهد. در پیاده‌سازی، شیب صفر و مخرج بسیار کوچک باید با قاعده عددی روشن مدیریت شوند.

\section{دو نوع بیش‌برازش}
بیش‌برازش متداول به افت تعمیم از داده آموزشی به داده آزمون اشاره دارد. بیش‌برازش ستبر در آموزش چندگامی نیز می‌تواند پس از مدتی به افت تدریجی درستی هم‌آورد آزمون منجر شود \cite{rice2020overfitting}. در مقابل، بیش‌برازش فاجعه‌بار معمولاً یک فروپاشی سریع همراه با شکاف میان حمله تک‌گامی و چندگامی است. یکسان‌گرفتن این پدیده‌ها سبب می‌شود راه‌حل‌هایی مانند توقف زودهنگام یا انتخاب نقطه کنترل بدون توجه به علت شکست به کار روند.

در این پژوهش، تمرکز بر نامعتبرشدن حل درونی تک‌گامی و نشانه‌های هندسی آن است. جلوگیری از این شکست، به‌تنهایی همه مسائل تعمیم، نامتوازن‌بودن دسته‌ها یا حساسیت به جابه‌جایی توزیع را حل نمی‌کند. تفکیک این اهداف کمک می‌کند نتایج تجربی با پرسش علمی مناسب سنجیده شوند.
